\section{Fundamentos de MoE: expertos, router y parámetros condicionales}

Solo después de definir con claridad los límites de recursos se puede definir Mixture of Experts (MoE, mezcla de expertos). Aquí, los «expertos» no son roles lingüísticos asignados manualmente, sino varias subredes que se aprenden; la «mezcla» tampoco significa ejecutar cada vez todas las subredes y promediarlas, sino que el router elige unas pocas rutas según la entrada. Esta idea es anterior al Transformer: en 1991, Adaptive Mixtures of Local Experts ya proponía que una gating network aprendiera la división del trabajo.

\lecturefigure{slide-04-adaptive-mixtures-history.jpg}{Hilo histórico desde Adaptive Mixtures of Local Experts hasta los modelos de lenguaje neuronales dispersos.}{Video oficial de Stanford Online 00:03:06--00:04:01.}

\begin{knowledgebox}{Glosario: primero, precisar el lenguaje del enrutamiento}
\begin{tabularx}{\textwidth}{L{0.22\textwidth}X}
\toprule
Término & Definición precisa en esta clase \\
\midrule
MoE & Un grupo de experts más un router; cada entrada activa solo unos pocos experts, por lo que los parámetros se usan de forma condicional. \\
expert & Suele ser un bloque independiente de parámetros de una feed-forward network (FFN); no equivale a una profesión interpretable por personas. \\
router / gating & Red ligera que calcula, a partir de la representación del token, las puntuaciones o probabilidades de los experts y decide la ruta de dispatch. \\
top-1 / top-2 & Cada token elige uno o dos experts de mayor puntuación, respectivamente; la segunda opción suele aumentar el cómputo y la comunicación. \\
sparse activation & La mayoría de los experts no se ejecuta para el token actual; es distinto de podar permanentemente los pesos a cero. \\
load balance & Hacer que la cantidad de tokens que recibe cada expert sea parecida, para facilitar multiplicaciones matriciales regulares por lotes. \\
\bottomrule
\end{tabularx}
\end{knowledgebox}

Sea la representación del token $x\in\mathbb{R}^{d}$, con $E$ experts y parámetros del router $W_r$. El gating estándar puede escribirse como:

\[
p_i(x)=\operatorname{softmax}(W_rx)_i,
\qquad
y(x)=\sum_{i\in\mathcal{S}(x)} g_i(x)E_i(x).
\]

Donde $p_i(x)$ es la probabilidad del router de que el token se asigne al $i$-ésimo expert; $\mathcal{S}(x)$ es el conjunto de experts elegido por top-$k$; $g_i(x)$ es el peso que, una vez hecha la elección, escala la salida; $E_i(x)$ es la salida de la FFN del $i$-ésimo expert; $y(x)$ es la salida de la capa dispersa. Solo se ejecutan los experts de $\mathcal{S}(x)$.

\lecturefigure{slide-05-moe-gating-equations.jpg}{Probabilidades de gating de MoE, selección de experts y salida ponderada.}{Video oficial de Stanford Online 00:04:02--00:04:25.}

\begin{importantbox}{La esencia de MoE}
MoE no consiste en «poner a votar varios modelos completos y promediarlos», sino en ejecutar un router muy pequeño para cada token y luego invocar unos pocos bloques de parámetros FFN. Permite que distintos tokens usen pesos distintos y, aun así, mantener un active compute similar.
\end{importantbox}

\teachervoice{La definición oral de Irwan Bello es importante: cada entrada recibe aproximadamente la misma cantidad de cómputo, pero se le aplican pesos distintos. Esta formulación distingue la sparsity de «simplemente reducir pasos de cómputo» y anticipa la discusión posterior sobre adaptive computation.}

Los primeros MoE dispersos lograron un éxito notable en traducción automática, pero su uso práctico exigía resolver tres problemas de sistemas. Primero, el enrutamiento hace que la cantidad de tokens por expert cambie dinámicamente; segundo, colocar los experts en distintos dispositivos genera comunicación; tercero, el router, la baja precisión y las fluctuaciones de carga amplifican la inestabilidad del entrenamiento. La contribución de Switch Transformer no fue proponer MoE por primera vez, sino reducir estos obstáculos a un esquema más simple, entrenable y escalable.

\lecturefigure{slide-06-moe-success-and-challenges.jpg}{Éxitos de los primeros MoE y sus obstáculos de complejidad, comunicación y estabilidad.}{Video oficial de Stanford Online 00:04:26--00:05:37.}

\begin{warningbox}{No tomar la expert specialization como un hecho a priori}
El router puede formar una división del trabajo por temas, idiomas o sintaxis, pero también puede aprender particiones numéricas más difíciles de interpretar. Los nombres de los expertos suelen ser resultado del análisis, no restricciones previas al entrenamiento; ver que cierto token va con más frecuencia a cierto expert no permite concluir directamente que ese expert posee una «habilidad» estable y causal.
\end{warningbox}

\subsection{Resumen de la sección}

MoE se compone de bloques de parámetros expert y un router; su núcleo es la selección condicional de parámetros por token. Amplía la capacidad disponible, pero introduce problemas de formas dinámicas, comunicación, estabilidad y balanceo de carga.
